% Generated draft: requires mathematical and source review.
\documentclass[11pt]{article}
\usepackage[margin=25mm]{geometry}
\usepackage{fontspec}
\setmainfont{Latin Modern Roman}
\usepackage{amsmath,amssymb,mathtools,unicode-math}
\setmathfont{Latin Modern Math}
\usepackage{enumitem,needspace,xurl,hyperref}
\setlength{\parindent}{0pt}
\setlength{\parskip}{5pt}
\setlength{\emergencystretch}{3em}
\newcommand{\N}{\mathbb{N}}
\newcommand{\Z}{\mathbb{Z}}
\newcommand{\Q}{\mathbb{Q}}
\newcommand{\R}{\mathbb{R}}
\newcommand{\C}{\mathbb{C}}
\newcommand{\F}{\mathbb{F}}
\newcommand{\A}{\mathbb{A}}
\newcommand{\PP}{\mathbb{P}}
\newcommand{\E}{\mathbb{E}}
\newcommand{\eps}{\varepsilon}
\newcommand{\dd}{\,\mathrm{d}}
\begin{document}
\section*{UnsolvedMath ID: 30001081}
\textbf{Source problem code:} catalogue entry
\subsection{Definitions and mathematical statement}
Let \(\rho:\R^3\to[0,\infty)\) be an unknown bounded probability density with compact support and zero center:
\[
\int_{\R^3}\rho(x)\,\dd x=1,\qquad
\int_{\R^3}x\rho(x)\,\dd x=0,\qquad \rho\in L^\infty(\R^3).
\]
For an orientation represented by a rotation \(R\), its two-dimensional marginal, in a fixed planar coordinate convention, is
\[
P_R\rho(z_1,z_2)=\int_\R\rho\bigl(R^{-1}(z_1,z_2,t)\bigr)\,\dd t.
\]
This formula explains the projection density rather than selecting a distribution of orientations.

Given finitely many such marginals in random unknown orientations, and \(n\) independent realizations of the resulting stochastic Radon transform, what can be inferred about \(\rho\)? Determine the information recoverable in the intended statistical model, distinguishing identification from a full observation law from inference with finitely many realizations. The source supplies no orientation law, noise or discretization model, loss function, or class beyond bounded centered compactly supported densities. It does not say whether the stated finite number of marginals is exactly \(n\), or each realization contains a collection. Consequently no unique consistency, reconstruction, confidence, or rate claim is specified. Orthogonal ambiguities mentioned in the assessment must be interpreted within the actual observation model rather than automatically asserting exact identification of \(\rho\).
\subsection{Short English statement}
Study what finitely many randomly oriented, unlabeled projection densities reveal about an unknown centered compactly supported density in three dimensions. The observation model and inferential target must be fixed, including which orthogonal symmetries are unobservable.
\subsection{Sources}
\begin{itemize}
\item Local UnsolvedMath registry, permanent ID \(30001081\), source code \texttt{catalogue entry}, ``Inference from Randomly Oriented Tomographic Marginals.'' Registry curation: Ulam AI / UnsolvedMath Contributors, CC BY 4.0; underlying sources retain their supplied provenance.
\item \emph{New Perspectives in Stochastic Geometry} (2009), Oberwolfach Reports, pp.\ 2655--2702, DOI \url{https://doi.org/10.4171/owr/2008/47}; supplied label ``OWR Report 47.''
\item Supplied \emph{On random tomography}: \url{https://doi.org/10.1214/08-AOS673}.
\item Supplied literature triage, \texttt{literature\_checked\_at}: 2026-08-21: labeled partially solved; says random tomography establishes identifiability modulo orthogonal symmetries and statistical procedures under models, while the broad finite-random-marginal inference question is not completely resolved. This is supplied dated curation, not independently verified; its model-dependent hypotheses are not present in the extraction.
\end{itemize}
\end{document}
